\documentclass[11pt]{article}
\usepackage[T1]{fontenc}
\usepackage[utf8]{inputenc}
\usepackage{lmodern,microtype}
\usepackage[margin=0.92in]{geometry}
\usepackage{amsmath,amssymb,amsthm,mathtools}
\usepackage{booktabs,tabularx,array,longtable}
\usepackage{enumitem}
\usepackage{xcolor}
\usepackage{hyperref}
\usepackage[nameinlink,capitalize,noabbrev]{cleveref}
\usepackage{fancyhdr,lastpage}

\definecolor{ink}{HTML}{17212B}
\definecolor{blue}{HTML}{295D84}
\definecolor{green}{HTML}{2C6749}
\definecolor{amber}{HTML}{895B14}
\hypersetup{colorlinks=true,linkcolor=blue,citecolor=blue,urlcolor=blue,
 pdftitle={Signature tori, fusion cycles, and the Hodge conjecture for p-ary covers},
 pdfauthor={Kyle Kistner}}
\pagestyle{fancy}\fancyhf{}
\fancyhead[L]{\small Signature tori and fusion cycles}
\fancyhead[R]{\small Phase II complete, 10 August 2026}
\fancyfoot[C]{\small \thepage\ of \pageref{LastPage}}
\renewcommand{\headrulewidth}{0.3pt}
\setlist[itemize]{leftmargin=1.45em,itemsep=0.22em,topsep=0.32em}
\setlist[enumerate]{leftmargin=1.7em,itemsep=0.22em,topsep=0.32em}
\renewcommand{\arraystretch}{1.14}

\newtheorem{theorem}{Theorem}[section]
\newtheorem{proposition}[theorem]{Proposition}
\newtheorem{lemma}[theorem]{Lemma}
\newtheorem{corollary}[theorem]{Corollary}
\theoremstyle{definition}
\newtheorem{definition}[theorem]{Definition}
\theoremstyle{remark}
\newtheorem{remark}[theorem]{Remark}
\crefname{theorem}{Theorem}{Theorems}
\crefname{proposition}{Proposition}{Propositions}
\crefname{lemma}{Lemma}{Lemmas}
\crefname{corollary}{Corollary}{Corollaries}
\crefname{definition}{Definition}{Definitions}
\crefname{remark}{Remark}{Remarks}

\newcommand{\Q}{\mathbb Q}
\newcommand{\Z}{\mathbb Z}
\newcommand{\F}{\mathbb F}
\newcommand{\C}{\mathbb C}
\newcommand{\R}{\mathbb R}
\newcommand{\PP}{\mathbb P}
\newcommand{\M}{\mathcal M}
\newcommand{\A}{\mathcal A}
\newcommand{\V}{\mathbb V}
\newcommand{\Mon}{\operatorname{Mon}}
\newcommand{\Hg}{\operatorname{Hg}}
\newcommand{\MT}{\operatorname{MT}}
\newcommand{\End}{\operatorname{End}}
\newcommand{\Hom}{\operatorname{Hom}}
\newcommand{\Gal}{\operatorname{Gal}}
\newcommand{\supp}{\operatorname{supp}}
\newcommand{\NS}{\operatorname{NS}}
\newcommand{\CH}{\operatorname{CH}}
\newcommand{\Sp}{\operatorname{Sp}}
\newcommand{\SU}{\operatorname{SU}}
\newcommand{\SL}{\operatorname{SL}}
\newcommand{\GL}{\operatorname{GL}}
\newcommand{\Res}{\operatorname{Res}}
\newcommand{\Pic}{\operatorname{Pic}}
\newcommand{\Alb}{\operatorname{Alb}}
\newcommand{\Jac}{\operatorname{Jac}}
\newcommand{\one}{\mathbf 1}
\newcommand{\eps}{\varepsilon}
\newcommand{\raw}{\mathrm{raw}}
\newcommand{\sat}{\mathrm{sat}}
\newcommand{\der}{\mathrm{der}}
\newcommand{\Nm}{\operatorname{Nm}}
\newcommand{\rk}{\operatorname{rk}}
\newcommand{\cl}{\operatorname{cl}}

\title{\vspace{-1.35em}\textbf{Signature Tori, Fusion Cycles, and the Hodge Conjecture\vspace{0.08em}\\for $p$-Ary Covers of the Projective Line}\vspace{0.15em}\\
\large Completion of the elementary-abelian Phase II programme}
\author{Kyle Kistner\\\small Ulam Capital\\\small \texttt{KyleJKistner@gmail.com}}
\date{10 August 2026}

\begin{document}
\maketitle
\vspace{-1.0em}

\begin{abstract}
Let $p$ be an odd prime and let a connected full family of $G=(\Z/p\Z)^m$-Galois covers of $\PP^1$ be encoded by its full-support branch code
\[
 C\subset \F_p^r\cap\one^\perp.
\]
The Phase I monodromy theorem determines the generic derived Hodge group and algebraizes every cross-character endomorphism defect.  We compute the missing center.  For a nonzero codeword $c$ put $s(c)=|\supp(c)|$ and
\[
 q_a(c)=\frac1p\sum_i[ac_i]_p,
 \qquad
 \delta_a(c)=s(c)-2q_a(c)
 \quad(a\in\F_p^\times).
\]
On the anti-invariant determinant lattice generated by the oriented character coordinates, the signature map is
\[
 \Sigma_C:\sum n_c\eps_c\longmapsto
 \left(\sum n_c\delta_a(c)\right)_{a\in\F_p^\times}.
\]
After quotienting by the Kummer matrix relations from Phase I, the character lattice of the generic central Hodge torus is exactly
\[
 X^*(T_C)=\operatorname{im}\Sigma_C,
\]
and the full generic Hodge group is the inverse image of $T_C$ in the saturated unitary centralizer.  Thus the complete generic Mumford--Tate group, including every internal rank-two corner, is determined by the code.

Over a splitting field the derived invariant algebra is generated by contractions and determinant tensors.  A determinant word is a Hodge tensor precisely when its exponent vector lies in $\ker\Sigma_C$.  For prime $p$, Aoki's theorem says that the concatenated residue tuple of every such relation is a union of opposite pairs.  We use those pairs to glue the associated cyclic quotient curves into a balanced compact-type admissible $\Z/p$-cover.  Smoothing gives a cyclic cover with simple branch tuple; Schoen's cycle theorem algebraizes its Weil determinant space.  Specialization through the compact-type relative Jacobian then produces the original mixed determinant tensor explicitly.  Consequently every rational Hodge class on every power of the very general Jacobian is algebraic.  More precisely, the Hodge ring is generated by divisor and Kummer-graph classes together with the compact-type fusion cycles indexed by $\ker\Sigma_C$.
\end{abstract}

\noindent\textbf{Proof status.}
Every Phase II dependency is either proved here, proved in the supplied Phase I manuscript, or quoted from a published theorem with its precise scope stated below.  The exact-arithmetic certificate independently checks the signature matrices, relation criterion, fusion rank identity, the first mixed $p=5$ class, the rank-two saturation criterion, and the split tower.  The manuscript has not undergone external specialist peer review.

\medskip
\noindent\textbf{Frozen scope.}
The theorem concerns very general Jacobians in full families of elementary abelian covers of odd prime exponent; the binary case is included separately at the end.  It determines the full generic Hodge/Mumford--Tate group and all rational Hodge classes on all powers.  Special-fibre jumps, composite-exponent non-simple balanced words, and arbitrary nonabelian deck groups are not asserted.

\tableofcontents

\section{Statement of the completed theorem}

Fix an odd prime $p$, put $G=(\Z/p\Z)^m$, and let
\begin{equation}\label{eq:branch}
 a_1,\ldots,a_r\in G\setminus\{0\},\qquad
 \sum_i a_i=0,\qquad
 \langle a_1,\ldots,a_r\rangle=G
\end{equation}
be connected branch data.  Let all ordered branch points vary over $B=\M_{0,r}$ and let
\[
 f:\mathcal X\longrightarrow B
\]
be the corresponding full family.  Evaluation identifies $G^\vee$ with a full-support code
\begin{equation}\label{eq:code}
 C=C_a\subset\F_p^r\cap\one^\perp.
\end{equation}
Let $A=JX_b$ for a very general $b\in B$.

For $x\in\F_p$, write $[x]_p\in\{0,1,\ldots,p-1\}$ for its least residue.  If $0\ne c=(c_i)\in C$, define
\begin{equation}\label{eq:signature-data}
 s(c)=|\supp(c)|,\qquad
 q_a(c)=\frac1p\sum_i[ac_i]_p,\qquad
 h(c)=s(c)-2,
\end{equation}
for $a\in\F_p^\times$, and
\begin{equation}\label{eq:delta}
 \delta_a(c)=s(c)-2q_a(c).
\end{equation}
The integer $h(c)$ is the rank over $K=\Q(\zeta_p)$ of the rational isotypic factor attached to the line $\F_pc$; the Chevalley--Weil multiplicities are
\begin{equation}\label{eq:CW}
 d_a(c)=s(c)-q_a(c)-1,\qquad
 d_{-a}(c)=q_a(c)-1,
\end{equation}
so $\delta_a(c)=d_a(c)-d_{-a}(c)$.

The determinant lattice and signature map will be defined intrinsically in \cref{sec:lattice}.  The following is the final result.

\begin{theorem}[Elementary-$p$ Phase II theorem]\label{thm:main}
Let $C\subset\F_p^r\cap\one^\perp$ be the full-support branch code of a connected full family of $(\Z/p)^m$-covers of $\PP^1$, where $p$ is odd, and let $A$ be the very general Jacobian.
\begin{enumerate}[label=\textup{(\roman*)}]
\item The complete connected generic Hodge group is
\[
 \Hg(A)=\pi^{-1}(T_C)\subset P_C,
\]
where $P_C$ is the saturated unitary centralizer from Phase I, $\pi:P_C\to D_C=P_C/P_C^{\der}$, and $T_C\subset D_C$ is the signature torus.  Its character lattice is
\[
 X^*(T_C)\cong\operatorname{im}\Sigma_C,
 \qquad
 \dim T_C=\rk\Sigma_C\le\frac{p-1}{2}.
\]
The generic Mumford--Tate group is $\MT(A)=\mathbb G_m\cdot\Hg(A)$.

\item Every generic Hodge endomorphism is determined by the code and is algebraic.  Cross-character corners are the Phase I Kummer correspondences.  Every residual internal standard--dual corner and every toric corner is a determinant relation in $\ker\Sigma_C$ and is realized by a fusion correspondence.

\item Over $\overline\Q$, the tensor invariants of $\Hg(A)$ are generated by contractions and determinant words whose exponent vectors lie in $\ker\Sigma_C$.  Contractions are represented by polarization and Poincar\'e divisor classes attached to the algebraic deck and Kummer homomorphisms.  Every permitted determinant word is represented by an explicit compact-type fusion cycle.

\item For every $q\ge1$, the rational Hodge ring of $A^q$ is generated by algebraic divisor/graph classes and fusion cycles.  In particular the rational Hodge conjecture holds for every power $A^q$.
\end{enumerate}
For $p=2$, the determinant-torus construction is unnecessary: the binary symplectic argument recalled in \cref{sec:HC} gives the same all-powers conclusion.
\end{theorem}

The remainder of the paper proves \cref{thm:main}.  The two new mechanisms are exact: the center is the Galois span of the determinant signature cocharacter, and the higher Hodge tensors are algebraized by compact-type fusion.

\section{Phase I input and the rational cyclic factors}\label{sec:phaseI}

We record only the Phase I result needed below.  Its proof, including the support reconstruction and four-point Kummer rigidity, is contained in the supplied companion manuscript \cite{KisI}.

For a nonzero line $\ell=\F_pc\subset C$, let
\[
 Y_c=X/\ker(c)\longrightarrow\PP^1
\]
be the cyclic quotient.  Since $p$ is prime and the quotient of $Y_c$ by $\Z/p$ is $\PP^1$, all of $H^1(Y_c,\Q)$ is primitive for the cyclotomic action.  It is a $K$-vector space of rank $h(c)=s(c)-2$ and depends only on $\ell$ as a rational Hodge structure.

\begin{theorem}[Phase I, exact derived group and algebraic blocks]\label{thm:phaseI}
After a finite \`etale base change, the following hold for the moving part of $R^1f_*\Q$.
\begin{enumerate}[label=\textup{(\alph*)}]
\item Every noncompact cyclic character factor has full simple algebraic monodromy projection.
\item Different supports separate.  On a common support of size at least five, a graph relation identifies only a character with itself or its inverse.  Every remaining relation is a four-point rank-two Kummer relation.
\item Every Kummer relation is induced, after finite base change, by an algebraic isomorphism of the associated cyclic quotient families.  Pulling and projecting its graph gives the corresponding Hodge homomorphism.
\item If $\A_{\sat}$ is generated by the deck projectors and all Kummer quotient graphs, and
\[
 P_C=C_{\Sp(H^1(A,\Q))}(\A_{\sat})^0,
\]
then
\[
 \Mon^0=P_C^{\der}=\Hg(A)^{\der}.
\]
All cross-character Hodge-homomorphism corners equal the corresponding corners of $\A_{\sat}$.
\end{enumerate}
\end{theorem}

The theorem was formulated for moving factors.  For odd prime exponent the omitted factors are completely understood.

\begin{lemma}[The only wholly toric factors have three branch points]\label{lem:toric-support3}
Let $0\ne c\in C$.  If $s(c)\ge4$, then $d_a(c)d_{-a}(c)>0$ for at least one $a\in\F_p^\times$.  Hence the rational factor $H^1(Y_c,\Q)$ has nontrivial derived monodromy.  A wholly toric factor has $s(c)=3$ and $h(c)=1$.
\end{lemma}

\begin{proof}
Suppose $s(c)\ge4$ and $d_a(c)d_{-a}(c)=0$ for every $a$.  The deformation space of polarized abelian varieties with the induced $K$-action has dimension
\[
 \frac12\sum_{a\in\F_p^\times}d_a(c)d_{-a}(c)=0.
\]
Thus the primitive Jacobian of the cyclic quotient family is constant.  Because $p$ is prime and the base quotient has genus zero, this primitive factor is the entire $JY_c$.  Torelli then makes the family $Y_c$ constant.  But the projection $\M_{0,r}\to\M_{0,s(c)}$ is dominant, and a fixed curve of genus at least two admits only finitely many faithful $\Z/p$-actions and quotient maps to $\PP^1$.  This contradicts $s(c)-3>0$.  Therefore some real place is noncompact.  If $s(c)=3$, the branch configuration is rigid and $h(c)=1$, so the factor is CM and toric.
\end{proof}

Thus \cref{thm:phaseI} supplies the complete derived group, while the support-three factors and the centers of the moving factors are exactly what the signature torus must recover.  From now on, \(\A_{\sat}\) denotes the Phase I algebra on the moving part, enlarged on the full \(H^1(A,\Q)\) by all deck idempotents and by every already constructed Kummer quotient graph.  We correspondingly redefine
\[
 P_C=C_{\Sp(H^1(A,\Q))}(\A_{\sat})^0.
\]
On a support-three factor this centralizer is its cyclotomic unitary torus and has trivial derived group.  Therefore the equality \(P_C^{\der}=\Hg(A)^{\der}\) remains valid on the full Hodge structure.  Phase II will cut down the center of this enlarged group and will produce the additional degree-two correspondences forced by that cutdown.

\section{The determinant lattice and the signature map}\label{sec:lattice}

Let
\[
 C^\circ=\{c\in C\setminus\{0\}:s(c)\ge3\}.
\]
We deliberately retain all nonzero scalar multiples: they are the different complex embeddings of the same rational cyclotomic factor.

\begin{definition}[Raw determinant lattice]\label{def:raw-lattice}
Set
\begin{equation}\label{eq:Eraw}
 E_C^{\raw}=
 \left(\bigoplus_{c\in C^\circ}\Z\eps_c\right)
 \Big/\langle\eps_{-c}+\eps_c:c\in C^\circ\rangle.
\end{equation}
The Galois group $\Gamma=\Gal(K/\Q)=\F_p^\times$ acts by
$a\eps_c=\eps_{ac}$.
\end{definition}

Over $\C$, $\eps_c$ is the determinant character of the complex factor $V_c$; polarization identifies $V_{-c}$ with $V_c^\vee$, giving $\eps_{-c}=-\eps_c$.

The Phase I Kummer matrix algebra imposes determinant equalities.  Let
\[
 R_{\mathrm{Kum}}\subset E_C^{\raw}
\]
be the lattice generated by $\eps_c-\eps_d$ when a Kummer matrix unit identifies $V_c$ with $V_d$, and by $\eps_c+\eps_d$ when it identifies $V_c$ with $V_d^\vee$.  Equivalently, it is the kernel of the natural determinant-character map from $E_C^{\raw}$ to the center quotient of the saturated centralizer.

\begin{proposition}[Character lattice of the saturated determinant quotient]\label{prop:D-lattice}
Let
\[
 D_C=P_C/P_C^{\der}.
\]
Then
\begin{equation}\label{eq:D-character}
 X^*(D_C)\cong E_C:=E_C^{\raw}/R_{\mathrm{Kum}}
\end{equation}
as $\Gamma$-lattices.
\end{proposition}

\begin{proof}
Work first over $\C$.  The exact block theorem writes the representation as a direct sum of standard and dual modules for the simple derived factors, with multiplicity spaces connected by the full Kummer matrix algebra.  Before those matrix identifications, the quotient by the derived group has one determinant character for every oriented coordinate $c$, and duality imposes $\eps_{-c}=-\eps_c$.  A standard-to-standard matrix unit identifies the two central scalars, while a standard-to-dual unit identifies one with the inverse of the other.  These are precisely the generators of $R_{\mathrm{Kum}}$.  No further relation occurs in the centralizer of the full matrix algebra.  The construction is Galois-equivariant, so descent gives \eqref{eq:D-character} over $\Q$.
\end{proof}

\begin{definition}[Signature map and relation lattice]\label{def:signature-map}
Define
\begin{equation}\label{eq:Sigma-raw}
 \widetilde\Sigma_C:E_C^{\raw}\longrightarrow\Z^{\Gamma},
 \qquad
 \widetilde\Sigma_C\!\left(\sum n_c\eps_c\right)
 =\left(\sum n_c\delta_a(c)\right)_{a\in\Gamma}.
\end{equation}
Put
\begin{equation}\label{eq:Lambda-raw}
 \widetilde\Lambda_C=\ker\widetilde\Sigma_C.
\end{equation}
\end{definition}

Because $\delta_a(-c)=-\delta_a(c)$, the map is well-defined.  Kummer-equivalent variations have equal determinant Hodge characters, with the sign prescribed by standard versus dual; hence
\begin{equation}\label{eq:R-in-Lambda}
 R_{\mathrm{Kum}}\subset\widetilde\Lambda_C.
\end{equation}
Consequently $\widetilde\Sigma_C$ descends to
\begin{equation}\label{eq:Sigma}
 \Sigma_C:E_C\longrightarrow\Z^{\Gamma},
 \qquad
 \Lambda_C:=\ker\Sigma_C=
 \widetilde\Lambda_C/R_{\mathrm{Kum}}.
\end{equation}
Since $\delta_{-a}(c)=-\delta_a(c)$,
\begin{equation}\label{eq:rank-bound}
 \rk\Sigma_C\le\frac{p-1}{2}.
\end{equation}

\begin{remark}[Exact computation]
Choose one representative of each sign orbit in $C^\circ$, form the integer matrix whose $c$th column is $(\delta_a(c))_{a\in\Gamma}$, and quotient its kernel by the explicit Kummer relations.  Smith normal form gives the integral relation lattice and the character lattice of $T_C$.  No period integration is required.
\end{remark}

\section{The exact generic central torus}\label{sec:center}

Let
\[
 h:S^1\longrightarrow\Hg(A)(\R)
\]
be the Hodge circle.  Write $\bar h$ for its projection to $D_C(\R)$.

\begin{lemma}[Determinant weight]\label{lem:det-weight}
For $c\in C^\circ$, the determinant character $\eps_c$ has weight
\[
 \langle\eps_c,\bar h\rangle=\delta_1(c).
\]
For the Galois conjugate $a\bar h$, the weight is $\delta_a(c)$.
\end{lemma}

\begin{proof}
On $V_c$, the Hodge circle acts by $z$ on the $d_1(c)$-dimensional $(1,0)$ part and by $z^{-1}$ on the $d_{-1}(c)$-dimensional $(0,1)$ part.  Its determinant is therefore $z^{d_1(c)-d_{-1}(c)}=z^{\delta_1(c)}$.  Applying $a\in\Gamma$ replaces $c$ by $ac$, giving $\delta_a(c)$.
\end{proof}

Let $T_C\subset D_C$ be the smallest $\Q$-subtorus containing $\bar h(S^1)$.  The next elementary torus lemma converts \cref{lem:det-weight} into an exact lattice formula.

\begin{lemma}[Galois-span criterion for a torus]\label{lem:torus-span}
Let $D$ be a $\Q$-torus, let $u:S^1\to D_\R$, and let $T\subset D$ be the smallest $\Q$-subtorus containing $u$.  A character $\chi\in X^*(D)$ is trivial on $T$ if and only if it has weight zero on every Galois conjugate of $u$.
\end{lemma}

\begin{proof}
Over $\overline\Q$, a subtorus is determined by its character annihilator.  The smallest $\Q$-subtorus containing $u$ is generated by the Galois orbit of the corresponding cocharacter.  A character annihilates that generated subtorus exactly when its pairing with every conjugate cocharacter is zero.
\end{proof}

\begin{theorem}[Signature-torus theorem]\label{thm:signature-torus}
The annihilator of $T_C$ in $X^*(D_C)=E_C$ is $\Lambda_C$.  Hence
\begin{equation}\label{eq:T-character}
 \boxed{X^*(T_C)\cong E_C/\Lambda_C
 \cong E_C^{\raw}/\widetilde\Lambda_C
 \cong\operatorname{im}\widetilde\Sigma_C.}
\end{equation}
In particular
\begin{equation}\label{eq:T-rank}
 \dim T_C=\rk\widetilde\Sigma_C\le\frac{p-1}{2}.
\end{equation}
\end{theorem}

\begin{proof}
By \cref{lem:det-weight,lem:torus-span}, a class represented by $e\in E_C^{\raw}$ annihilates $T_C$ precisely when
\[
 \langle e,a\bar h\rangle=0\quad\text{for every }a\in\Gamma,
\]
which is exactly $\widetilde\Sigma_C(e)=0$.  Since $R_{\mathrm{Kum}}\subset\widetilde\Lambda_C$, the annihilator in $E_C$ is $\widetilde\Lambda_C/R_{\mathrm{Kum}}=\Lambda_C$.  Quotienting character lattices gives \eqref{eq:T-character}; the quotient is torsion-free because it is identified with the image of a homomorphism between free lattices.  The rank assertion follows from \eqref{eq:rank-bound}.
\end{proof}

\begin{theorem}[Full generic Hodge and Mumford--Tate groups]\label{thm:full-Hg}
With $\pi:P_C\to D_C=P_C/P_C^{\der}$,
\begin{equation}\label{eq:Hg-preimage}
 \boxed{\Hg(A)=\pi^{-1}(T_C).}
\end{equation}
Moreover
\begin{equation}\label{eq:MT}
 \MT(A)=\mathbb G_m\cdot\Hg(A).
\end{equation}
\end{theorem}

\begin{proof}
Phase I gives $\Hg(A)^{\der}=P_C^{\der}$.  The image of $\Hg(A)$ in $D_C$ is, by the Tannakian definition of the generic Hodge group, the smallest $\Q$-subtorus containing the projected Hodge circle, namely $T_C$.  Therefore
\[
 \Hg(A)\subseteq\pi^{-1}(T_C)
\]
and both connected reductive groups have dimension
$\dim P_C^{\der}+\dim T_C$.  They are equal.  For a polarized weight-one Hodge structure, adjoining homotheties to the Hodge group gives the Mumford--Tate group, proving \eqref{eq:MT}.
\end{proof}

The formerly unresolved rank-two corner is now explicit.

\begin{corollary}[Internal disk criterion]\label{cor:disk-corner}
Let $c$ be a four-point rank-two coordinate not already saturated by a Kummer matrix block.  The derived-invariant standard--dual homomorphism is a Hodge tensor if and only if
\begin{equation}\label{eq:zero-signature-disk}
 \delta_a(c)=0\qquad\text{for every }a\in\F_p^\times.
\end{equation}
When this occurs, the four residues of $c$ are two opposite pairs.  A double transposition carries $c$ to $-c$, and the corresponding Phase I Kummer graph algebraizes the internal corner.
\end{corollary}

\begin{proof}
The standard--dual homomorphism has central character $2\eps_c$.  Since $X^*(T_C)$ is torsion-free, it is trivial on $T_C$ precisely when $\eps_c\in\widetilde\Lambda_C$, which is \eqref{eq:zero-signature-disk}.  This says that the four-entry residue tuple is balanced in the sense of \cref{def:balanced} below.  Aoki's prime-exponent theorem, recalled in \cref{thm:Aoki}, makes it a union of opposite pairs.  Exchanging the two pairs by the appropriate double transposition sends the tuple to its negative, and Phase I Kummer rigidity supplies the graph correspondence.
\end{proof}

\section{Derived invariants and determinant words}\label{sec:invariants}

We now pass from endomorphisms to all tensor invariants.  Work temporarily over an algebraic closure.  The exact derived-group theorem decomposes $H^1(A)$ into standard and dual modules for a product of special linear groups, with algebraic matrix units on every Kummer multiplicity space.  The support-three factors have rank one and hence no derived group.

\begin{proposition}[First fundamental theorem in the present representation]\label{prop:FFT}
The invariants of $\Hg(A)^{\der}$ in the tensor algebra of $H^1(A^q)$ are generated by:
\begin{enumerate}[label=\textup{(\alph*)}]
\item contractions between a standard module and its dual; and
\item determinant tensors $\det(V_c)=\bigwedge^{h(c)}V_c$ and their duals.
\end{enumerate}
Arbitrary choices in the Kummer multiplicity spaces are obtained from the algebraic matrix units of $\A_{\sat}$.
\end{proposition}

\begin{proof}
After scalar extension, each nontrivial derived factor is an $\SL_h$ acting on copies of its standard representation and its dual.  The first fundamental theorem for $\SL_h$ says that invariants in mixed tensor powers are generated by evaluation pairings and the two volume tensors in $\bigwedge^hV$ and $\bigwedge^hV^\vee$ \cite{Weyl}.  The Kummer matrix algebra is full on each multiplicity space by Phase I, so its algebraic matrix units transport a chosen contraction or volume tensor to every copy.
\end{proof}

The contractions are already algebraic: polarizations and Poincar\'e divisor classes attached to the deck and Kummer homomorphisms give the required degree-two classes; the corresponding homomorphisms may equivalently be represented by their algebraic graphs.  The determinant tensors carry the entire residual central character.

\begin{definition}[Effective determinant word]\label{def:word}
An effective determinant word is a finite ordered list
\[
 \mathbf c=(c_1,\ldots,c_N),\qquad c_j\in C^\circ.
\]
Its exponent is
\[
 [\mathbf c]=\sum_{j=1}^N\eps_{c_j}\in E_C^{\raw},
\]
and its determinant space over $\C$ is
\begin{equation}\label{eq:det-space-complex}
 \mathcal D(\mathbf c)_\C=
 \bigoplus_{a\in\F_p^\times}
 \bigotimes_{j=1}^N\det(V_{ac_j}).
\end{equation}
The Galois-stable rational form is denoted $\mathcal D(\mathbf c)$.
\end{definition}

A negative coefficient in a lattice relation is made effective by replacing $c$ with $-c$; polarization identifies the resulting determinant with the dual, up to an algebraic Tate factor.

\begin{proposition}[Hodge criterion for a determinant word]\label{prop:Hodge-word}
The following are equivalent:
\begin{enumerate}[label=\textup{(\roman*)}]
\item $\mathcal D(\mathbf c)$ is a Hodge substructure of type
$\bigl(H/2,H/2\bigr)$, where $H=\sum_jh(c_j)$;
\item $[\mathbf c]\in\widetilde\Lambda_C$;
\item for every $a\in\F_p^\times$,
\begin{equation}\label{eq:word-signature}
 \sum_{j=1}^N\delta_a(c_j)=0.
\end{equation}
\end{enumerate}
In that case $H$ is even.
\end{proposition}

\begin{proof}
At the embedding $a$, the determinant $\det(V_{ac_j})$ has Hodge bidegree
$(d_a(c_j),d_{-a}(c_j))$.  The tensor product has equal bidegrees precisely when the difference \eqref{eq:word-signature} vanishes.  This is exactly the definition of $\widetilde\Lambda_C$.  Equal bidegrees sum to $H/2$, so $H$ is even.
\end{proof}

\begin{theorem}[Complete tensor reduction]\label{thm:tensor-reduction}
Every rational Hodge tensor on every power of $A$ is a linear combination of tensor contractions and determinant spaces $\mathcal D(\mathbf c)$ with $[\mathbf c]\in\widetilde\Lambda_C$.  Conversely every such contraction and determinant space consists of Hodge classes.
\end{theorem}

\begin{proof}
By \cref{prop:FFT}, every derived invariant is generated by contractions and determinant tensors.  The full Hodge group is the inverse image of $T_C$ by \cref{thm:full-Hg}.  A determinant monomial survives the center exactly when its character is trivial on $T_C$, equivalently when its exponent lies in $\widetilde\Lambda_C$ by \cref{thm:signature-torus}.  Contractions have trivial determinant character.  This proves both directions after scalar extension; taking Galois invariants gives the rational statement.
\end{proof}

Thus Phase II reduces the Hodge conjecture to one exact geometric statement: every balanced determinant word must be algebraic.

\section{Prime balance and Aoki's theorem}\label{sec:Aoki}

For a list $\mathbf c=(c_j)$, concatenate all nonzero entries of the codewords, with repetitions, to obtain
\[
 \beta(\mathbf c)=(\beta_1,\ldots,\beta_S)
 \in(\F_p^\times)^S,
 \qquad
 S=\sum_js(c_j).
\]

\begin{definition}[Balanced and simple tuples]\label{def:balanced}
A tuple $\beta\in(\F_p^\times)^S$ is \emph{balanced} if
\begin{equation}\label{eq:balanced}
 \sum_{i=1}^S[a\beta_i]_p=\frac{pS}{2}
 \qquad\text{for every }a\in\F_p^\times.
\end{equation}
It is \emph{simple} if, after a permutation, it has the form
\[
 (x_1,-x_1,\ldots,x_{S/2},-x_{S/2}).
\]
\end{definition}

\begin{lemma}[Relations are balanced tuples]\label{lem:relation-balanced}
For an effective determinant word $\mathbf c$,
\[
 [\mathbf c]\in\widetilde\Lambda_C
 \quad\Longleftrightarrow\quad
 \beta(\mathbf c)\text{ is balanced}.
\]
\end{lemma}

\begin{proof}
For every $a$,
\[
 \sum_j\delta_a(c_j)
 =S-\frac2p\sum_i[a\beta_i]_p.
\]
The left side vanishes exactly when \eqref{eq:balanced} holds.
\end{proof}

The arithmetic input is the following theorem of Aoki, in precisely the form quoted and used by Schoen \cite{Aoki}\cite[Corollary~1.9 and Aoki's theorem]{Schoen}.

\begin{theorem}[Aoki]\label{thm:Aoki}
If $p$ is prime, every balanced tuple in $(\F_p^\times)^S$ is simple.  Equivalently, in a balanced tuple each residue $x$ occurs with the same multiplicity as $-x$.
\end{theorem}

\begin{remark}
The prime hypothesis is the exact boundary used here.  Aoki's theorem also covers $m=4$ and certain length-dependent composite cases, but balanced non-simple tuples exist for general composite exponent.  Phase II makes no composite-exponent claim.
\end{remark}

Combining \cref{lem:relation-balanced,thm:Aoki}, every determinant relation carries a concrete opposite-residue pairing.  The next section turns such a pairing into an algebraic cycle.

\section{Compact-type fusion}\label{sec:fusion}

Fix an effective word
\[
 \mathbf c=(c_1,\ldots,c_N),\qquad
 [\mathbf c]\in\widetilde\Lambda_C.
\]
Let
\[
 Y_j=Y_{c_j}=X/\ker(c_j)\longrightarrow\PP^1
\]
be the corresponding cyclic quotient curves, allowing repetitions.  We identify each quotient group $G/\ker(c_j)$ with the same abstract group $\F_p$ through $c_j$; under this identification its local monodromy labels are exactly the nonzero entries of $c_j$.  Aoki's theorem pairs every branch occurrence carrying residue $x$ with an occurrence carrying residue $-x$.

\subsection{The pairing graph}

Make a graph $\mathcal G_{\mathbf c}$ whose vertices are $1,\ldots,N$ and whose edges are the chosen opposite pairs; loops are allowed.  For each connected component choose a spanning tree using cross-vertex edges.  An edge in the spanning tree will become a node; all other opposite pairs remain smooth branch points.

\begin{lemma}[Fusion bookkeeping]\label{lem:fusion-rank}
Let a connected component have vertices $J$ and $k=|J|$.  If its spanning tree uses $k-1$ opposite pairs as nodes, then the smoothed cyclic cover has
\begin{equation}\label{eq:fused-branch}
 s_{\mathrm{fus}}=\sum_{j\in J}s(c_j)-2(k-1)
\end{equation}
smooth branch points and $K$-rank
\begin{equation}\label{eq:fused-rank}
 h_{\mathrm{fus}}=s_{\mathrm{fus}}-2
 =\sum_{j\in J}\bigl(s(c_j)-2\bigr)
 =\sum_{j\in J}h(c_j).
\end{equation}
The remaining branch tuple is simple.
\end{lemma}

\begin{proof}
Each tree edge removes two branch markings and replaces them by one node on each of the two source components.  This gives \eqref{eq:fused-branch}.  Since a tree has $k-1$ edges,
\[
 s_{\mathrm{fus}}-2
 =\sum_js(c_j)-2k
 =\sum_j(s(c_j)-2).
\]
Every opposite pair not used by the tree survives as two smooth branch points, so the remaining tuple is still a union of opposite pairs.
\end{proof}

\subsection{The admissible cyclic cover}

Take one copy $P_j\cong\PP^1$ for every vertex.  On $P_j$ mark the branch points of $Y_j$ with their residues.  For every spanning-tree edge pairing residues $x$ and $-x$, glue the two marked points on the bases and glue the unique totally ramified points above them.  The resulting source and target have tree dual graphs.

\begin{proposition}[Balanced compact-type admissible cover]\label{prop:admissible}
The construction gives a connected balanced admissible $\Z/p$-cover
\[
 Y_0\longrightarrow P_0
\]
with the following properties.
\begin{enumerate}[label=\textup{(\alph*)}]
\item $P_0$ is a nodal genus-zero curve with tree dual graph.
\item $Y_0$ has the same tree dual graph and is of compact type.
\item The normalization of $Y_0$ is the disjoint union of the curves $Y_j$.
\item After a finite base change, the cover occurs as the special fibre of a one-parameter family whose generic fibre is a smooth cyclic $p$-cover of $\PP^1$ with the simple branch tuple of \cref{lem:fusion-rank}.
\end{enumerate}
\end{proposition}

\begin{proof}
At a glued branch point the two local inertia characters are $\alpha$ and $-\alpha$.  Since $\alpha\ne0$ and $p$ is prime, both points are totally ramified and there is one point above each.  Locally the balanced node has the form
\[
 uv=t,\qquad \xi\eta=s,\qquad u=\xi^p,\quad v=\eta^p,\quad t=s^p,
\]
with the deck generator acting by $(\xi,\eta)\mapsto(\zeta_p^{\alpha}\xi,\zeta_p^{-\alpha}\eta)$.  Thus the two tangent characters are inverse.  The base and source dual graphs are the chosen trees, proving (a)--(c).  The displayed local model smooths every node after the finite base change $t=s^p$.  The target tree has unobstructed deformations; for a tame finite group the cotangent complex of $B(\Z/p)$ is zero, so the cover contributes no additional global obstruction, and the balanced local smoothing parameters may be chosen simultaneously.  Equivalently, the balanced twisted-cover stack is smooth and gives the normalization of the admissible-cover compactification \cite{ACV}.  This proves (d).  The branch residues left on the generic smooth base are exactly those not used as nodes.
\end{proof}

For a compact-type family, the identity component of the relative Picard functor is an abelian scheme.  Hence a smoothing in \cref{prop:admissible} gives an abelian scheme $\mathcal J\to\operatorname{Spec}R$ over a discrete valuation ring with
\begin{equation}\label{eq:special-Jacobian}
 \mathcal J_0\cong\prod_{j\in J}JY_j,
 \qquad
 \mathcal J_\eta\cong JY_\eta.
\end{equation}
The cyclotomic action extends over $R$.

\subsection{Schoen cycles and specialization}

We use the following published result.  In Schoen's notation, the indicated Hodge structure is the cyclotomic determinant, or Weil, subspace.

\begin{theorem}[Schoen's simple-tuple cycle theorem]\label{thm:Schoen}
Let $Y\to\PP^1$ be a cyclic cover of degree $p$ with simple branch tuple, and let
\[
 H=\dim_KH^1(Y,\Q).
\]
If $H$ is even, then the rational determinant Hodge subspace
\[
 \mathcal W(Y)=\bigwedge_K^H H^1(Y,\Q)
 \subset H^H(JY,\Q)(H/2)
\]
is generated by classes of codimension $H/2$ algebraic cycles on $JY$.
\end{theorem}

\begin{proof}[Published input]
Schoen's Theorem~2.0 constructs the cycles on the required self-product of the curve when the branch tuple is simple, and Corollary~3.1 pushes them to the primitive abelian factor of the Jacobian \cite{Schoen}.  For a prime cyclic cover of $\PP^1$, every nontrivial character is primitive and the invariant part of $H^1$ is zero, so that primitive factor is the whole Jacobian.
\end{proof}

\begin{lemma}[Specialization of the determinant space]\label{lem:specialization}
In the compact-type smoothing \eqref{eq:special-Jacobian}, the specialization isomorphism carries
\[
 \mathcal W(Y_\eta)
\]
onto
\begin{equation}\label{eq:det-direct-sum}
 \bigwedge_K^{\sum_jh(c_j)}
 \left(\bigoplus_jH^1(Y_j,\Q)\right)
 \cong
 \bigotimes_j\bigwedge_K^{h(c_j)}H^1(Y_j,\Q).
\end{equation}
If $\mathcal W(Y_\eta)$ is generated by algebraic cycles, then the right-hand determinant space is generated by algebraic cycles on $\prod_jJY_j$.
\end{lemma}

\begin{proof}
Smooth proper base change for the abelian scheme $\mathcal J/R$ gives a $K$-linear identification of the generic and special $H^1$.  The special fibre decomposition \eqref{eq:special-Jacobian} identifies the latter with the displayed direct sum.  The canonical determinant identity gives \eqref{eq:det-direct-sum}.

After a finite extension of the fraction field, choose algebraic cycles spanning $\mathcal W(Y_\eta)$ and take their closures in a proper model of the relevant power of $\mathcal J$.  The specialization map on Chow groups is compatible with the cohomological specialization isomorphism \cite{Fulton,SGA6}.  Their specializations therefore span the right-hand side of \eqref{eq:det-direct-sum}.
\end{proof}

\begin{theorem}[Fusion algebraization theorem]\label{thm:fusion}
For every effective determinant word $\mathbf c$ with
$[\mathbf c]\in\widetilde\Lambda_C$, the rational Hodge space
$\mathcal D(\mathbf c)$ is generated by algebraic cycles on
\[
 \prod_{j=1}^N JY_{c_j}.
\]
After pull--push through the quotient maps $X\to Y_{c_j}$ and the rational group-algebra projectors, the same Hodge space is generated by algebraic cycles on a power of $A=JX$.  These are the \emph{fusion cycles} attached to $\mathbf c$ and the chosen opposite-residue pairing.
\end{theorem}

\begin{proof}
By \cref{lem:relation-balanced,thm:Aoki}, the concatenated branch tuple admits an opposite-residue pairing.  Apply the pairing-graph construction component by component.  For one connected component, \cref{prop:admissible} supplies a compact-type smoothing whose generic cyclic cover has a simple branch tuple and whose $K$-rank is the sum of the component ranks by \cref{lem:fusion-rank}.  This rank is even because the remaining fused branch tuple is simple, hence has even length.  Schoen's theorem algebraizes the generic determinant space; \cref{lem:specialization} specializes it to the tensor product of the determinant spaces of the component curves.  Taking external products over pairing-graph components gives $\mathcal D(\mathbf c)$.

For each quotient map $q_c:X\to Y_c$, the maps $q_c^*:JY_c\to JX$ and $(q_c)_*:JX\to JY_c$ are algebraic and satisfy $(q_c)_*q_c^*=[\deg(q_c)]=[p^{m-1}]$.  The rational group-algebra idempotent isolates the required line factor.  Products, pullbacks, pushforwards, and rational linear combinations of the specialized cycles therefore give algebraic cycles on a power of $A$ with exactly the desired Betti realization.
\end{proof}

\begin{remark}[Constructivity]
A fusion cycle is determined by finite data: a lattice relation, an opposite-residue matching, a spanning forest, Schoen's explicit simple-tuple cycles on the fused cover, and specialization.  Different matchings can give different algebraic representatives of the same rational Hodge subspace.
\end{remark}

\section{The Hodge conjecture on every power}\label{sec:HC}

\begin{theorem}[All-powers Hodge theorem]\label{thm:all-powers}
Let $p$ be an odd prime, let $A$ be the very general Jacobian in a connected full family of $(\Z/p)^m$-covers of $\PP^1$, and let $q\ge1$.  Every rational Hodge class on $A^q$ is algebraic.  More precisely, the rational Hodge classes on all powers of $A$ are generated, in the tensor category under products, pullbacks, pushforwards, diagonals, and permutations, by:
\begin{enumerate}[label=\textup{(\alph*)}]
\item polarization and Poincar\'e divisor classes;
\item Poincar\'e divisor classes attached to the deck-projector and Kummer homomorphisms from Phase I (equivalently, their algebraic graph correspondences); and
\item fusion cycles for determinant relations in $\widetilde\Lambda_C$.
\end{enumerate}
\end{theorem}

\begin{proof}
By \cref{thm:tensor-reduction}, every Hodge tensor is a linear combination of contractions and determinant spaces associated with $\widetilde\Lambda_C$.  The contractions are represented by the divisor and graph classes in (a) and (b).  Every determinant space is algebraic by \cref{thm:fusion}.  The cohomology of an abelian variety is the exterior algebra on $H^1$; its Hodge classes are obtained from the tensor invariants by alternating projectors, which are algebraic combinations of permutations.  Therefore every Hodge class on $A^q$ is represented by an algebraic cycle.
\end{proof}

\begin{corollary}[Finite code-theoretic generating data]\label{cor:finite-generation}
Choose an integral basis of $\widetilde\Lambda_C$ by Smith normal form and orient each basis vector as an effective word.  Together with the Phase I divisor/graph classes, the corresponding fusion classes and their duals generate the tensor category of rational Hodge classes on all powers of $A$.
\end{corollary}

\begin{proof}
Products add determinant exponents, while replacing $c$ by $-c$ and using polarization gives the dual relation.  Hence an integral lattice basis generates every determinant character relation.  Matrix units and permutations supply all multiplicity choices.
\end{proof}

\begin{corollary}[Complete generic endomorphism theorem]\label{cor:endomorphisms}
Every element of $\End_{\mathrm{HS}}(H^1(A,\Q))$ is algebraic.  Its matrix corners are computable from $\Sigma_C$ and the Phase I Kummer blocks.  In particular
\[
 \End^0(A)=\End_{\Hg(A)}H^1(A,\Q)
\]
is the algebra generated by the deck action, Phase I quotient graphs, and degree-two fusion correspondences.
\end{corollary}

\begin{proof}
Endomorphisms are the degree-two part of the tensor-invariant theorem.  Cross-character derived invariants are already Kummer graphs.  Every remaining central condition is a determinant relation of total cohomological degree two and is algebraized by fusion.  Full faithfulness of the weight-one Hodge realization identifies Hodge endomorphisms with rational endomorphisms of $A$.
\end{proof}

For $p=2$, every nonzero character is self-dual and Phase I gives a product of full symplectic groups.  The first fundamental theorem for the symplectic group leaves only contractions, so divisor classes generate the Hodge ring on all powers.  This completes \cref{thm:main} for every prime.

\section{Examples and diagnostics}\label{sec:examples}

\subsection{The split elementary-\texorpdfstring{$p$}{p} tower}

For
\[
 C_{\mathrm{split}}=
 \{(a,b,-a,-b):a,b\in\F_p\},
\]
every nonzero codeword is already a union of opposite pairs.  Hence
\[
 \delta_u(c)=0
 \qquad(c\in C_{\mathrm{split}},\ u\in\F_p^\times),
\]
so
\begin{equation}\label{eq:split-torus-zero}
 T_{C_{\mathrm{split}}}=1.
\end{equation}
This explains, rather than merely verifies, the semisimple Hodge group and divisor-generated Hodge ring obtained in Phase I.  All determinant characters are relations, but their classes are already generated by the global Kummer matrix algebra and $\SL_2$ contractions.

\subsection{The first mixed \texorpdfstring{$p=5$}{p=5} fusion class}

Let $p=5$ and take the length-eight code generated by
\begin{align*}
 c&=(1,1,3,0,0,0,0,0),\\
 d&=(0,0,0,1,2,4,4,4).
\end{align*}
Then
\[
 (\delta_a(c))_{a=1}^4=(1,1,-1,-1),
 \qquad
 (\delta_a(d))_{a=1}^4=(-1,-1,1,1),
\]
so
\begin{equation}\label{eq:mixed-relation}
 \eps_c+\eps_d\in\widetilde\Lambda_C.
\end{equation}
Neither determinant line is Hodge by itself.  Their product is a codimension-two Hodge space.

The concatenated tuple is
\[
 (1,1,3,1,2,4,4,4).
\]
Pair the residue $3$ on the first quotient with the residue $2=-3$ on the second and use it as the unique node.  The fused smooth branch tuple is
\[
 (1,1,1,4,4,4),
\]
which is simple and has $K$-rank
\[
 (3-2)+(5-2)=4.
\]
Schoen's codimension-two cycles on the fused Jacobian specialize to the mixed determinant class in \eqref{eq:mixed-relation}.  This is the first example in the programme where the algebraic class is genuinely mixed across a rigid CM factor and a moving rank-three factor.

\subsection{Rank and defect diagnostics}

The signature matrix distinguishes three phenomena:
\begin{enumerate}
\item $\rk\Sigma_C=0$: the connected Hodge group is semisimple modulo homotheties; every determinant tensor is Hodge.
\item $0<\rk\Sigma_C<(p-1)/2$: the code has nontrivial Weil relations but a positive-dimensional central torus.
\item $\rk\Sigma_C=(p-1)/2$: the central torus has maximal possible rank; exceptional determinant tensors are exactly the nontrivial integral dependencies among the codeword signature columns.
\end{enumerate}
The relation lattice is therefore both the exact Hodge-tensor classifier and an effective computational invariant of the branch code.

\section{What Phase II completes}\label{sec:completion}

For elementary abelian covers of odd prime exponent, together with the separate binary case, the programme is now closed at the generic Hodge-theoretic level:
\[
 \boxed{
 \begin{array}{c}
 \text{$p$-ary branch code}\\[2pt]
 \Downarrow\\[-1pt]
 \text{exact derived monodromy and Kummer blocks (Phase I)}\\[2pt]
 \Downarrow\\[-1pt]
 \text{signature matrix and full central Hodge torus}\\[2pt]
 \Downarrow\\[-1pt]
 \text{all determinant/Weil Hodge relations}\\[2pt]
 \Downarrow\\[-1pt]
 \text{compact-type fusion and Schoen cycles}\\[2pt]
 \Downarrow\\[-1pt]
 \text{rational Hodge conjecture on every power.}
 \end{array}}
\]

The proof separates the two logically different sources of Hodge tensors.  Derived-group contractions are algebraic divisors and graph classes.  Central-torus relations are exactly balanced residue words, and prime balance makes each of them accessible to Schoen's simple-tuple construction after fusion.  No unidentified generic tensor remains.

The boundary of the theorem is exact.  Composite exponent permits balanced non-simple tuples, so Aoki's prime simplification fails and Schoen's theorem no longer automatically algebraizes every determinant relation.  Special fibres can have additional Mumford--Tate drops.  Nonabelian covers require a different representation-theoretic replacement for the cyclotomic determinant lattice.  None of those questions is imported into the result proved here.

\appendix
\section{Certificate ledger}

The attached script
\begin{quote}\small
\path{phase_II_certificate.py}
\end{quote}
performs exact integer and rational arithmetic checks of the following interfaces:
\begin{enumerate}
\item construction of the signature matrix from a full-support $p$-ary code;
\item Smith-normal-form $\Z$-bases of the exact relation lattices for representative codes, and the equality between signature relations and balanced concatenated tuples;
\item opposite-residue multiplicities for every tested relation;
\item the compact-type fusion rank formula component by component;
\item the mixed $p=5$ relation and its fused six-point tuple;
\item the zero-signature disk/Kummer criterion through $p=19$;
\item vanishing of the signature torus for the split elementary-$p$ family through $p=31$.
\end{enumerate}
The finite computation is a regression certificate.  Aoki's theorem, admissible-cover smoothing, specialization of cycles, Schoen's theorem, the central-torus argument, and invariant theory are written mathematical inputs, not computational assertions.

\begin{thebibliography}{99}

\bibitem{ACV}
D. Abramovich, A. Corti, and A. Vistoli,
\emph{Twisted bundles and admissible covers},
Comm. Algebra 31 (2003), 3547--3618.

\bibitem{Aoki}
N. Aoki,
\emph{On some arithmetic problems related to the Hodge cycles on the Fermat varieties},
Math. Ann. 266 (1983), 23--54;
erratum, Math. Ann. 267 (1984), 572.

\bibitem{BL}
C. Birkenhake and H. Lange,
\emph{Complex Abelian Varieties}, second edition,
Grundlehren Math. Wiss. 302, Springer, 2004.

\bibitem{DeligneCycles}
P. Deligne,
\emph{Hodge cycles on abelian varieties},
in: Hodge Cycles, Motives, and Shimura Varieties,
Lecture Notes in Math. 900, Springer, 1982, 9--100.

\bibitem{Fulton}
W. Fulton,
\emph{Intersection Theory}, second edition,
Ergebnisse Math. Grenzgeb. 2, Springer, 1998.

\bibitem{KisI}
K. Kistner,
\emph{Abelian branch codes, exact monodromy blocks, and algebraic Mumford--Tate defects},
Phase I working manuscript, 2026; supplied with this package.

\bibitem{Schoen}
C. Schoen,
\emph{Hodge classes on self-products of a variety with an automorphism},
Compositio Math. 65 (1988), 3--32.

\bibitem{SGA6}
P. Berthelot, A. Grothendieck, L. Illusie, et al.,
\emph{Th\'eorie des intersections et th\'eor\`eme de Riemann--Roch},
SGA 6, Lecture Notes in Math. 225, Springer, 1971.

\bibitem{Weyl}
H. Weyl,
\emph{The Classical Groups: Their Invariants and Representations},
Princeton University Press, 1939.

\end{thebibliography}

\end{document}
